%=============================================================
%  Module I.4 / L06 — Zeeman Effect
%  Series I > Module I.4 > Lecture L06
%  Duration: 90 minutes  |  All tiers (T1–T5)
%=============================================================
\renewcommand{\lectureCode}{I.4 / L06}

\section*{L06 — Zeeman Effect}
\label{sec:L06}
\addcontentsline{toc}{section}{L06 — Zeeman Effect}

%--- Lecture header
\begin{tcolorbox}[lecturebox, title={Module I.4 / L06 — Metadata}]
\begin{tabular}{@{}ll@{}}
\textbf{Series:}    & I — Introductory Quantum Mechanics \\
\textbf{Module:}    & I.4 — Paradigm Physical Systems \\
\textbf{Lecture:}   & L06 — Zeeman Effect \\
\textbf{Duration:}  & 90 minutes \\
\textbf{Tiers:}     & T1 (HS) · T2 (BegUG) · T3 (AdvUG) · T4 (MSc) · T5 (PhD) \\
\textbf{Preceding:} & I.4 / L05 — Minimal Coupling and Landau Levels \\
\textbf{Following:} & I.4 / L07 — Stark Effect \\
\textbf{Prerequisites:} & Modules I.1–I.3 complete \\
\end{tabular}

\medskip
\textbf{Key concepts:} Normal Zeeman triplet, Landé g-factor, anomalous Zeeman, Paschen–Back limit, Breit–Rabi
\end{tcolorbox}

%------------------------------------------------------------
\subsection*{L06.0 — Learning Objectives (per tier)}
\label{sec:L06.0}

\tiertag{tier1col}{Tier 1 — HS}\quad
Identify the main physical phenomenon studied in this lecture; describe it in
non-mathematical terms; state one real-world application.

\tiertag{tier2col}{Tier 2 — BegUG}\quad
Apply the key equations to compute numerical results; verify consistency with
known limits; translate between physical and mathematical descriptions.

\tiertag{tier3col}{Tier 3 — AdvUG}\quad
Derive central results from first principles; prove key identities; identify
the mathematical structure underlying the physical phenomenon.

\tiertag{tier4col}{Tier 4 — MSc}\quad
Analyse formal extensions; connect to symmetry principles; generalise to
related systems; identify the role of this lecture in the broader programme.

\tiertag{tier5col}{Tier 5 — PhD}\quad
Establish rigorous mathematical foundations; identify open research problems;
connect to modern literature; critique standard textbook treatments.

%--- Lecture content -----------------------------------------

\subsection*{L06.1 — Normal Zeeman Effect}
Orbital magnetic moment $\hat{\boldsymbol\mu}_L = -\mu_B\hat{\mathbf{L}}/\hbar$;
perturbation $\hat{H}'=\mu_B B\hat{L}_z/\hbar$.
First-order shift: $E^{(1)}_{nlm} = \mu_B m B$; triplet splitting $\Delta E = \mu_B B$.

\subsection*{L06.2 — Anomalous Zeeman and Land\'e g-Factor}
\begin{tcolorbox}[lecturebox, title={Landé g-Factor}]
\begin{equation}
  g_J = 1 + \frac{j(j+1)+s(s+1)-l(l+1)}{2j(j+1)},
  \quad E^{(1)}_{njm_j} = g_J\mu_B m_j B.
  \label{eq:L06.1}
\end{equation}
\end{tcolorbox}

\subsection*{L06.3 — Paschen–Back and Breit–Rabi}
Strong-field limit ($\mu_B B \gg$ SO splitting): good quantum numbers switch
$j,m_j \to m_l,m_s$; energies linear in $m_l + 2m_s$.
Breit–Rabi diagram: exact interpolation between weak and strong field for
two-state mixing.


%------------------------------------------------------------
\subsection*{L06.C — Connections}
\label{sec:L06.C}
\textbf{Backward:} I.4/L05 — Minimal Coupling and Landau Levels and Modules I.1–I.3.
\textbf{Forward:} I.4/L07 — Stark Effect builds directly on the formalism developed here.
\textbf{Cross-module:} Series~II modules extend these results to many-body and
advanced settings; Series~III applies them to molecular electronic structure.

%=============================================================
%  ASSESSMENT BUNDLE — L06
%=============================================================
\newpage
\section*{L06 — Assessment Artifact Bundle}

\begin{tcolorbox}[ugconceptbox, title={SCQU · L06 · Tiers 1–3 · 15–20 min}]
\textit{Ten simple concept questions (mix MC/TF/short) targeting the core results
of Zeeman Effect. See artifact file \texttt{L06_artifacts.tex} for full questions.}

\textbf{Rubric:} 2 pts each (1.5 correct + 0.5 justification). Total: 20 pts.
\end{tcolorbox}

\begin{tcolorbox}[ugconceptbox, title={CCQU · L06 · Tiers 2–3 · 25–35 min}]
\textit{Ten complex concept questions: ``explain why,'' ``what would change if,''
``identify the flaw.'' Full questions in \texttt{L06_artifacts.tex}.}

\textbf{Rubric:} 2 pts each (1 key point + 1 argument). Total: 20 pts.
\end{tcolorbox}

\begin{tcolorbox}[ugprobbox, title={SPSU · L06 · Tiers 1–3 · 60–90 min}]
\textit{Ten problems (Part A: mechanics, Part B: interpretation, Part C: translation).
Full problems in \texttt{L06_artifacts.tex}.}

\textbf{Rubric:} Result 60\%, Method 30\%, Interpretation 10\%.
\end{tcolorbox}

\begin{tcolorbox}[ugprobbox, title={CPSU · L06 · Tier 3 · 3–6 hrs (4-layer)}]
\textit{Five multi-part derivation problems with four-layer scaffolding.
Layers 1–4 in \texttt{L06_ex01\_problem.tex}, \texttt{hints}, \texttt{adv\_hints}, \texttt{solution}.}

\textbf{Rubric:} Assumptions 15\%, Proof structure 45\%, Algebra 25\%, Insight 15\%.
\end{tcolorbox}

\begin{tcolorbox}[ugprobbox, title={SPJU + CPJU · L06 · Tiers 2–3}]
\textit{Simple project (2–6 hrs): structured computational/analytical task.
Complex project (8–15 hrs): open-ended synthesis.
Full specs in \texttt{L06_artifacts.tex}.}
\end{tcolorbox}

\begin{tcolorbox}[gradconceptbox, title={SCQG + CCQG · L06 · Tiers 4–5}]
\textit{Ten simple + ten complex concept questions at graduate level (rigour, domain
conditions, named theorems). Full questions in \texttt{L06_artifacts.tex}.}
\end{tcolorbox}

\begin{tcolorbox}[gradprobbox, title={SPSG + CPSG · L06 · Tiers 4–5}]
\textit{Ten simple + five complex graduate problems.
CPSG with four-layer format. See \texttt{L06_artifacts.tex}.}
\end{tcolorbox}

\begin{tcolorbox}[gradprobbox, title={SPJG (×5) + CPJG (×5) · L06 · Tiers 4–5}]
\textit{Five simple (3–8 hrs each) and five complex (10–20 hrs each) graduate projects.
Full specs in \texttt{L06_artifacts.tex}.}
\end{tcolorbox}

\begin{tcolorbox}[researchbox, title={WRQ (×5) + ORQ (×5) · L06 · Tiers 4–5}]
\textit{Five well-defined research questions (5–10 hrs each) and five open-ended
research questions (10–20 hrs each). Full prompts in \texttt{L06_artifacts.tex}.}
\end{tcolorbox}
